Collaborative \emph{machine learning} (ML) seeks to build ML models of higher quality by training on more data owned by multiple parties \cite{simcollaborative,yang2019federated}. For example, a hospital can improve its prediction of disease progression by training on data collected from more and diversified patients from other hospitals \cite{share-health}. Likewise, a real-estate firm can improve its prediction of demand and price by training on data from others \cite{share-real-estate}.  However, parties have two main concerns that discourage data sharing and participation in collaborative ML: (a) whether they benefit from the collaboration and (b) privacy.

Concern (a) arises as each party would expect to cover the significant cost they incur to collect and share data (e.g., the risk of losing their competitive edge).
Some existing works \cite{simcollaborative,tay2021}, among other data valuation methods \footnote{Data valuation methods study how much data is worth. As explained in~\citet{sim2022data}, a party's data is first valued independently using a performance metric (e.g., see Def.~\ref{def:bs} later) and then relative to the data contributed by others (e.g., Shapley value (Sec.~\ref{sec:incentives})). The latter value is helpful to (i) model interpretability and (ii) deciding how much to compensate the data owners \emph{fairly}.}, have recognized that parties require incentives to collaborate, such as a guaranteed \emph{fair} higher reward from contributing more valuable data than the others, an \emph{individually rational} higher reward from collaboration than in solitude, and a higher total reward value (\emph{group welfare}) whenever possible. Often, parties share and are rewarded with information (e.g., gradients \cite{xu2021gradient} or parameters \cite{simcollaborative} of parametric ML models) computed from the shared data. However, these incentive schemes expose parties to privacy risks.

On the other hand, some \emph{federated learning} (FL) works  \cite{mcmahan2017seminalFL} have addressed the privacy concern (b) and satisfied strict data protection laws (e.g., European Union's General Data Protection Regulation) by enforcing \emph{differential privacy} (DP) \cite{abadi2016deep,mcmahan2017learning} during the collaboration. Each party injects noise before sharing information to ensure that its shared information would not significantly alter a knowledgeable collaborating party's or mediator's belief about whether a datum was input to the algorithm. Injecting more noise leads to a stronger DP guarantee.
As raised in \citet{Yu21}, adding DP can invalidate game-theoretic properties and hence affect participation. For example, in the next paragraph, we will see that adding DP may lead to the collaboration being perceived as unfair and a lower group welfare.
However, to the best of our knowledge (and as discussed in Sec.~\ref{sec:related} and Fig.~\ref{fig:related-venn}), there are no works that address both concerns, i.e., ensure the \emph{fairness}, \emph{individual rationality}, and \emph{group welfare} incentives (see Sec.~\ref{sec:incentives}), alongside privacy.
Thus, we aim to fill in this gap and design an incentive scheme with privacy by addressing the following questions:
